\documentclass{article}

\usepackage[utf8]{inputenc}
\usepackage{graphicx}
\usepackage{geometry, float}
\usepackage[style=authoryear, backend=biber]{biblatex}
\usepackage{mhchem}
\usepackage{amsmath}

\geometry{ a4paper, top=2cm, bottom=2.5cm, left=3cm, right=3cm, heightrounded}
\renewcommand{\baselinestretch}{1.15}
\setlength{\parindent}{0pt}
\setlength{\parskip}{0.8em}

\addbibresource{references.bib}

\newcommand{\figref}[1]{Figure~\ref{#1}}

\title{Dissertation Proposal}
\author{Jakub Watson-Jones}
\date{10 February 2026}

\begin{document} \maketitle 

\section{Context}
Maths education is crucial to our every day lives. A good understanding of maths and logic allows informed decision making, creative problem solving, strategic planing etc. The UK government recognizes education as a instrument of economic growth. The government also uses education as a soft power in international politics~(\cite{UK_IES_2026}). The net education related exports generate about £32.3 billion~(\cite{UKrevenue}). The UK deepens relationships and trust with other countries using education and by educating 59 foreign world leaders~(\cite{HEPI_SPI}).  

However, today we often find a single teacher teaching a large number of students this makes the teacher's attention a limited resource. Each student learns at their own pace resulting in the teacher having to focus on slower learners. A single teacher in charge can also easily miss individual gaps in students' understanding. Today, we have algorithms able to generate personalized feedback for each individual. Software like Sparx Maths, Century tech, Mathletics Time Table Rock Stars, utilise algorithms that generate personalized feedback, search for gaps in knowledge, and generate reports for teachers, schools and parents.




Education is crucial to our every day lives. Traditional education suffers from limited teacher time for each student. High number of students for each teacher results in either one of two downsides. Either each students receives the same attention, or teachers focus on the slower students attempting to get each student at the same level of 


\section{Aim}





\section{Planning}
\subsection{Work Plan}
\subsection{Work Plan Explanation}
\subsection{Risk}

\newpage
\printbibliography{}
\end{document}